\subsection{Dimension and chains of graded ideals}

In this issue, we will denote by dimgr(H) the supremum of the lengths of chains of graded prime ideals of H; likewise, if p is a graded prime ideal of H, we will denote by htgr(p) the supremum of the lengths of chains of graded prime ideals of H for which p is the largest element. If p is a graded prime ideal of H, we have \(p \cap S = \varnothing\); otherwise, indeed p would contain an element whose degree-0 component equals 1, hence would contain 1 since it is graded. The map \(p \mapsto S^{-1}p\) from the set of graded prime ideals of H to the set of prime ideals of \(S^{-1}H\) is therefore injective and increasing (II, § 2, no. 5, prop. 11); consequently, in view of § 1, no. 3, prop. 6 and the corollary to prop. 7, we have:

PROPOSITION 2. --- a) We have \(\dim \mathrm{gr}(\mathrm{H}) \leqslant \dim (\mathrm{S}^{-1}\mathrm{H}) \leqslant \dim (\mathrm{H})\) .

\begin{enumerate}
\def\labelenumi{\alph{enumi})}
\setcounter{enumi}{1}
\tightlist
\item
  For every graded prime ideal \(\mathfrak{p}\) from H, we have \(\operatorname{htgr}(\mathfrak{p}) \leqslant \operatorname{ht}(\mathrm{S}^{-1}\mathfrak{p}) = \operatorname{ht}(\mathfrak{p})\) .
\end{enumerate}

For every ideal a of H, denote \(a^{gr}\) the largest graded ideal contained in a; we have \(a^{gr} = \sum_{n} (a \cap H_n)\) .

Lemma 1. --- a) If p is a prime ideal of H, \(\mathfrak{p}^{\mathrm{gr}}\) is a prime ideal.

\begin{enumerate}
\def\labelenumi{\alph{enumi})}
\setcounter{enumi}{1}
\item
  Every maximal element of the set of graded ideals of H distinct from H is a maximal ideal of H that contains \(H_{\geq1}\) .
\item
  Every minimal prime ideal of H is graded.
\item
  This follows from III, § 1, no. 4, prop. 4.
\item
  Let m be a graded ideal of H distinct from H. Then we have
\end{enumerate}

\[
\mathfrak {m} \subset (\mathfrak {m} \cap \mathrm{H} _ {0}) + \mathrm{H} _ {\geqslant 1} \neq \mathrm{H}.
\]

If m is maximal, then we have \(m = m_{0} + H_{\geq 1}\) , where \(m_{0}\) is a maximal ideal of \(H_{0}\) , whence b).

\begin{enumerate}
\def\labelenumi{\alph{enumi})}
\setcounter{enumi}{2}
\tightlist
\item
  Let p be a minimal prime ideal of H. Since \(p^{gr}\) is prime by a) and contained in p, we have \(p = p^{gr}\) whence c).
\end{enumerate}

Lemma 2.---Let p and q be prime ideals of H such that q ⊂ p and q ≠ p. If q \(^{gr}\) = p \(^{gr}\) then q is graded, p is not, and ht(p/q) = 1.

Remark 1.--- Let us take up the notations of example 1 of no. 1. Lemma 2 implies that, if two irreducible subvarieties Y and Z of \(\mathbf{C}^{n+1}\) have the same projecting cone and if \(Z \subset Y\) and \(Z \neq Y\) , then Y is the projecting cone of Z, and Z has codimension 1 in Y. *

Replacing H by \(\mathrm{H} / \mathfrak{q}^{\mathrm{gr}}\) we reduce to the case where \(\mathfrak{q}^{\mathrm{gr}} = \{0\}\). Then H is integral (lemma 1, a)), \(\mathfrak{p}^{\mathrm{gr}} = 0\) , and the task is to prove that \(\operatorname{ht}(\mathfrak{p}) \leqslant 1: this will indeed imply that  \operatorname{ht}(\mathfrak{q}) = 0\) , hence that \(\mathfrak{q} = \{0\}\) . Since \(\mathfrak{p}^{\mathrm{gr}} = \{0\}\) , one has \(\mathfrak{p} \cap \mathrm{H}_n = \{0\}\) for every \(n\) , and \(\mathfrak{p}\) is disjoint from the multiplicative subset \(\mathrm{T} = \bigcup_{n} (\mathrm{H}_n - \{0\})\) . The ring \(\mathrm{H}_{\mathfrak{p}}\) is therefore isomorphic to a ring of fractions of \(\mathrm{T}^{-1}\mathrm{H}\) , and we therefore have

\[
\mathrm{ht} (\mathfrak {p}) = \dim (\mathrm{H} _ {\mathfrak {p}}) \leqslant \dim (\mathrm{T} ^ {- 1} \mathrm{H})
\]

(§ 1, no. 3, props. 6 and 7). But, by Lemma 4 of V, § 1, no. 8, T \(^{-1}\) H is a field or is isomorphic to a ring K{[}X, X \(^{-1}\) {]}, where K is a field; we therefore have dim(T \(^{-1}\) H) ≤ 1 and ht(p) ≤ 1, which is what we wanted to prove.

PROPOSITION 3. --- Let p be a prime ideal of H. If p ≠ p \(^{gr}\) , we have ht(p \(^{gr}\) ) = ht(p) - 1. By Lemma 1, a), the ideal p \(^{gr}\) is prime and contained in p, hence ht(p \(^{gr}\) ) ≤ ht(p) - 1. The proposition being trivial when ht(p \(^{gr}\) ) = +∞, we may assume ht(p \(^{gr}\) ) \textless{} + ∞. Let us prove the inequality ht(p) ≤ ht(p \(^{gr}\) ) + 1 by induction on ht(p \(^{gr}\) ). It suffices to prove that, for every prime ideal q contained in p and distinct from p, one has ht(q) ≤ ht(p \(^{gr}\) ). We distinguish two cases according as q \(^{gr}\) ≠ p \(^{gr}\) or that q \(^{gr}\) = p \(^{gr}\) . If q \(^{gr}\) ≠ p \(^{gr}\) , then we have ht(q \(^{gr}\) ) \textless{} ht(p \(^{gr}\) ); we have

\[
\mathrm{ht} (\mathfrak {q}) \leqslant \mathrm{ht} (\mathfrak {q} ^ {\mathrm{gr}}) + 1,
\]

by the induction hypothesis if q ≠ q \(^{gr}\) and trivially if q = q \(^{gr}\) ; consequently, we have ht(q) ≤ ht(q \(^{gr}\) ) + 1 ≤ ht(p \(^{gr}\) ), which is what we wanted to prove. If q \(^{gr}\) = p \(^{gr}\) , then we have q = q \(^{gr}\) by Lemma 2, hence ht(q) = ht(q \(^{gr}\) ) ≤ ht(p \(^{gr}\) ), whence again the desired conclusion.

THEOREM 1. --- Assume H is Noetherian.

\begin{enumerate}
\def\labelenumi{\alph{enumi})}
\item
  Every chain of graded prime ideals of H, saturated as a chain of graded prime ideals, is saturated as a chain of prime ideals.
\item
  For every graded prime ideal p of H, one has htgr(p) = ht(S \(^{-1}\) p) = ht(p).
\item
  We have \(\dim \mathrm{gr}(\mathbf{H}) = \dim (\mathbf{S}^{-1}\mathbf{H}) = \dim (\mathbf{H}).\)
\end{enumerate}

To prove \(a\) ), it suffices to show that, if \(\mathfrak{p}\) and \(\mathfrak{q}\) are two distinct graded prime ideals of H such that \(\mathfrak{q} \subset \mathfrak{p}\) and every graded prime ideal between \(\mathfrak{q}\) and \(\mathfrak{p}\) is equal to \(\mathfrak{p}\) or to \(\mathfrak{q}\) , then \(\mathrm{ht}(\mathfrak{p}/\mathfrak{q}) = 1\) . By dividing by \(\mathfrak{q}\) , one reduces to the case where \(\mathfrak{q} = \{0\}\) . It therefore remains to prove that, if H is integral, and if \(\mathfrak{p}\) is an ideal of H, minimal among the graded prime ideals \(\neq \{0\}\) , one has \(\mathrm{ht}(\mathfrak{p}) = 1\) . Now let \(a\) be a nonzero homogeneous element of \(\mathfrak{p}\) , and let r be a prime ideal of H such that \(a \in \mathfrak{r} \subset \mathfrak{p}\) and minimal for these properties (II, § 2, no. 6, lemma 2). Since \(\mathfrak{r}^{\mathrm{gr}}\) is prime (lemma 1, \(a\) ) and nonzero, we have \(\mathfrak{r}^{\mathrm{gr}} = \mathfrak{p}\) , hence \(\mathfrak{p} = \mathfrak{r}\) . Since H is integral and Noetherian, \(\mathfrak{p}\) has height 1 (§ 3, no. 1, prop. 1), whence \(a\) ).

Let us prove b). Let p be a graded prime ideal of H. We have

\[
\operatorname{htgr} (\mathfrak {p}) \leqslant \operatorname{ht} \left(\mathrm{S} ^ {- 1} \mathfrak {p}\right) \leqslant \operatorname{ht} (\mathfrak {p})
\]

(prop. 2); let us prove the inequality ht(p) ≤ htgr(p) by induction on htgr(p). If htgr(p) = 0, p is minimal among graded prime ideals, hence minimal (lemma 1, c)), and we have ht(p) = 0. Suppose that htgr(p) \textgreater{} 0 and let us prove the inequality ht(q) ≤ htgr(p) - 1 for every prime ideal q contained in p and distinct from p. Let us distinguish two cases. If q is graded, we conclude by the induction hypothesis. If q is not graded, then we have q \(^{gr}\) ≠ p, hence ht(q \(^{gr}\) ) ≤ htgr(q \(^{gr}\) ) by the induction hypothesis, whence ht(q) ≤ htgr(q \(^{gr}\) ) + 1 by prop. 3; it remains to prove the inequality htgr(q \(^{gr}\) ) ≤ htgr(p) - 2; but if we had htgr(q \(^{gr}\) ) = htgr(p) - 1, the chain q \(^{gr}\) ⊂ p would be saturated by a), which is not the case since q \(^{gr}\) ≠ q ≠ p.

Let us finally prove c). We have dimgr(H) ≤ dim(S⁻¹H) ≤ dim(H) (prop. 2), and it remains to prove dim(H) ≤ dimgr(H), or equivalently ht(p) ≤ dimgr(H) for every prime ideal p of H. Let p therefore be a prime ideal of H. If p is graded, we have ht(p) = htgr(p) ≤ dimgr(H). If p is not graded, we have ht(p) = htgr(p\(^{\mathrm{gr}}\)) + 1 by prop. 3; let m be a maximal graded ideal of H containing p\(^{\mathrm{gr}}\); by lemma 1, b), m is maximal, hence distinct from p\(^{\mathrm{gr}}\), and we have htgr(p\(^{\mathrm{gr}}\)) + 1 ≤ htgr(m) ≤ dimgr(H), whence again ht(p) ≤ dimgr(H). This completes the proof.

Remark 2. --- There exist non-Noetherian graded rings H such that dimgr(H) \textless{} dim(H) (p.~99, exercise 1).

